\begin{abstract}
\noindent\textbf{Summary:} \rev{GSEA is a standard approach for pathway interpretation, yet Python ecosystems lack a high-performance implementation aligned with the fgseaMultilevel rare-event estimator target, especially for trajectory-aware rolling-window analysis. PyFgsea achieves numerical alignment with the R fgseaMultilevel reference for normalized enrichment scores (NES; Pearson $r>0.999$) and significantly reduces runtime compared to existing Python tools. Its stateful rolling-window engine further reduces repeated trajectory-analysis overhead, yielding approximately 1.9-fold end-to-end wall-time speedup in a conservative stress test and up to 7.47-fold component-level acceleration in a narrower 100-window benchmark.}

\medskip
\noindent\textbf{Availability and implementation:} Source code is available at \url{https://github.com/shayuanxukuang/pyfgsea} and via PyPI (\texttt{pip install pyfgsea}). \rev{An archival snapshot of the code and benchmark data is available on Zenodo (DOI: \href{https://doi.org/10.5281/zenodo.19446446}{10.5281/zenodo.19446446}).}

\medskip
\noindent\textbf{Contact:} \href{mailto:shih@kust.edu.cn}{shih@kust.edu.cn}

\medskip
\noindent\textbf{Supplementary information:} Supplementary data are available at Bioinformatics online.

\medskip
\noindent\textbf{Keywords:} Gene set enrichment analysis; multilevel sampling; rare-event estimation; single-cell trajectories; Rust; PyO3
\end{abstract}
